\documentclass[11pt,a4paper]{article}
\usepackage{fontspec}
\usepackage[vietnamese]{babel}
\usepackage{geometry}
\geometry{a4paper, top=14mm, bottom=16mm, left=18mm, right=18mm, includehead, headheight=26pt, headsep=6mm}
\usepackage{amsmath,amssymb}
\usepackage{graphicx}
\usepackage{xcolor}
\usepackage{listings}
\usepackage{booktabs}
\usepackage{tcolorbox}
\usepackage{fancyhdr}
\usepackage{hyperref}

% Hệ màu Brand Identity Asian School
\definecolor{AsianBlue}{RGB}{2, 103, 177}       % #0267b1
\definecolor{AsianDarkBlue}{RGB}{26, 84, 147}   % #1a5493
\definecolor{AsianLightBlue}{RGB}{238, 246, 252} 

\pagestyle{fancy}
\fancyhf{}
\renewcommand{\headrulewidth}{1pt}
\renewcommand{\headrule}{\hbox to\headwidth{\color{AsianBlue}\leaders\hrule height \headrulewidth\hfill}}
\lhead{\color{AsianDarkBlue}\small\textbf{Trường Quốc tế Á Châu} -- Asian School}
\rhead{\color{AsianDarkBlue}\small\textbf{Lý thuyết Tuần 01} -- Olympic 11}
\rfoot{\color{gray}\small Trang \thepage}

% Cấu hình listings
\definecolor{codegreen}{rgb}{0,0.55,0}
\definecolor{codegray}{rgb}{0.5,0.5,0.5}
\definecolor{backcolour}{RGB}{248, 250, 252}

\lstdefinestyle{asianstyle}{
    backgroundcolor=\color{backcolour},   
    commentstyle=\color{codegreen}\itshape,
    keywordstyle=\color{AsianDarkBlue}\bfseries,
    numberstyle=\tiny\color{codegray},
    stringstyle=\color{red!70!black},
    basicstyle=\ttfamily\footnotesize,
    breakatwhitespace=false,         
    breaklines=true,                 
    keepspaces=true,                 
    numbers=left,
    numbersep=6pt,
    tabsize=4,
    frame=single,
    rulecolor=\color{AsianBlue!40},
    showstringspaces=false
}
\lstset{style=asianstyle}

\begin{document}

% BANNER ĐẦU TRANG
\noindent
\begin{minipage}[c]{0.38\textwidth}
    \includegraphics[height=1.35cm]{logo_asian_school.png}
\end{minipage}
\hfill
\begin{minipage}[c]{0.58\textwidth}
    \raggedleft
    {\fontsize{9.5pt}{11.5pt}\selectfont\bfseries\color{AsianDarkBlue}TRƯỜNG QUỐC TẾ Á CHÂU -- ASIAN SCHOOL}\\[2pt]
    {\fontsize{8.5pt}{10.5pt}\selectfont\color{gray}TỔ TIN HỌC -- BỘ PHẬN BỒI DƯỠNG NĂNG KHIẾU}\\[2pt]
    {\fontsize{8.5pt}{10pt}\selectfont\color{AsianBlue}\textbf{ĐỘI TUYỂN OLYMPIC TIN HỌC 11 -- NĂM HỌC 2026--2027}}
\end{minipage}

\vspace{0.1cm}
{\color{AsianBlue}\hrule height 1.5pt}
\vspace{0.15cm}

\begin{center}
    {\Large\bfseries\color{AsianDarkBlue}TÀI LIỆU LÝ THUYẾT TUẦN 01}\\[3pt]
    {\normalsize\textbf{\color{AsianBlue}NHẬP MÔN C++ THI ĐẤU \& TƯ DUY NỀN TẢNG TỪ GỐC}}\\[2pt]
    \textit{\small Bản chất biến nhớ RAM, kiểu dữ liệu, toán tử, rẽ nhánh và vòng lặp}
\end{center}

\vspace{-0.1cm}
\begin{tcolorbox}[colback=AsianLightBlue,colframe=AsianBlue,arc=1.5mm,title=\textbf{\footnotesize MỤC TIÊU BÀI HỌC CỐT LÕI}]
\footnotesize
\begin{enumerate}
    \setlength{\itemsep}{1pt}
    \item \textbf{Định hướng \& Giải thưởng:} Nắm vững thể lệ kỳ thi Olympic 11 TP.HCM, cơ cấu giải và phổ điểm mục tiêu.
    \item \textbf{Tư duy biến nhớ \& Máy tính:} Hiểu bản chất vùng nhớ RAM, tại sao C++ bắt buộc khai báo kiểu; khắc phục triệt để bẫy tràn số nguyên (\textit{integer overflow}) và bẫy chia nguyên.
    \item \textbf{Làm chủ cú pháp C++:} Lệnh nhập/xuất \texttt{cin}, \texttt{cout}, cấu trúc rẽ nhánh \texttt{if/else}, vòng lặp \texttt{for} kèm bảng chạy từng bước (\textit{dry run}).
    \item \textbf{Thực hành Code::Blocks:} Làm chủ phím tắt \texttt{F9}, cấu hình cờ \texttt{-std=c++17}, đọc xuất tệp \texttt{.INP} / \texttt{.OUT}.
\end{enumerate}
\end{tcolorbox}

\vspace{-0.2cm}
\section*{\color{AsianDarkBlue}\normalsize 1. Thông tin Kỳ thi \& Phổ điểm chuẩn Olympic 11 TP.HCM}
\vspace{-0.1cm}
\begin{itemize}
    \setlength{\itemsep}{1.5pt}
    \small
    \item \textbf{Thời gian \& Quy cách:} Thi vào đầu tháng 04/2027 (120 phút). Đề thi gồm \textbf{3 bài toán} (thang 20 điểm: 6đ -- 7đ -- 7đ). Chấm tự động qua tệp \texttt{<TENBAI>.INP} và \texttt{<TENBAI>.OUT}.
    \item \textbf{Tỷ lệ trao giải:} Chiếm 50\% -- 60\% thí sinh dự thi:
\end{itemize}

\vspace{-0.2cm}
\begin{center}
\footnotesize
\begin{tabular}{@{}lcp{8.5cm}@{}}
\toprule
\textbf{Hạng giải} & \textbf{Điểm chuẩn} & \textbf{Chiến lược làm bài phòng thi} \\
\midrule
\textbf{Huy chương Vàng} & 16.0 -- 20.0 đ & Trọn vẹn Bài 1 (6đ) + Bài 2 (7đ) + Ăn 3--4đ subtask Bài 3 \\
\textbf{Huy chương Bạc} & 13.0 -- 15.5 đ & Trọn vẹn Bài 1 (6đ) + Ăn 5--6đ Bài 2 + Ăn 1--2đ subtask Bài 3 \\
\textbf{Huy chương Đồng} & 10.0 -- 12.5 đ & Trọn vẹn Bài 1 (6đ) + Ăn 4--5đ subtask cơ bản Bài 2 \\
\textbf{Khuyến khích} & 8.0 -- 9.5 đ & Trọn vẹn Bài 1 (6đ) + Ăn 2--3đ subtask vét cạn Bài 2 \\
\bottomrule
\end{tabular}
\end{center}

\vspace{-0.1cm}
\begin{tcolorbox}[colback=white,colframe=AsianDarkBlue,arc=1mm,boxrule=0.5pt]
\footnotesize\textbf{Mục tiêu chiến lược của Đội tuyển Asian School:} Vùng điểm \textbf{12.0 -- 15.0 điểm} để nắm chắc \textbf{Huy chương Bạc / Đồng}. Thuật toán vét cạn vẫn ăn được 30\% -- 60\% số điểm của mỗi bài!
\end{tcolorbox}

\vspace{-0.2cm}
\section*{\color{AsianDarkBlue}\normalsize 2. Vì sao C++ là vũ khí số 1 trong Lập trình thi đấu?}
\vspace{-0.1cm}
\begin{itemize}
    \setlength{\itemsep}{1.5pt}
    \small
    \item \textbf{Tốc độ thực thi vượt trội:} C++ biên dịch trực tiếp sang mã máy nhị phân, xử lý được $\approx 10^8$ phép tính/giây (nhanh hơn Python 10 -- 100 lần), triệt tiêu nguy cơ quá thời gian (\textbf{TLE}).
    \item \textbf{Thư viện chuẩn STL cực mạnh:} Tích hợp sẵn \texttt{std::sort}, \texttt{std::vector}, tìm kiếm nhị phân tối ưu.
    \item \textbf{Môi trường Code::Blocks:} IDE chuẩn mực được cài đặt sẵn trên máy thi chính thức của Sở GD\&ĐT.
\end{itemize}

\newpage

\section*{\color{AsianDarkBlue}\normalsize 3. Bản chất Biến nhớ \& Tư duy Máy tính từ con số 0}

\subsection*{\color{AsianBlue}\small a) Biến nhớ trong RAM là gì?}
Khi máy tính chạy chương trình, mọi dữ liệu được lưu trên bộ nhớ tạm thời \textbf{RAM}. RAM gồm hàng tỷ ô nhớ.\\
\textbf{Hình ảnh ẩn dụ:} Hãy tưởng tượng một \textbf{Biến} giống như một \textbf{Chiếc hộp có nhãn}:
\begin{itemize}
    \setlength{\itemsep}{1pt}
    \small
    \item \textbf{Kiểu dữ liệu:} Hình dạng/kích cỡ chiếc hộp (chỉ đựng số nguyên, chỉ đựng số thực, hay ký tự).
    \item \textbf{Tên biến:} Nhãn dán bên ngoài hộp để lập trình viên gọi tên (\texttt{tuoi}, \texttt{tong}, \texttt{dem}).
    \item \textbf{Giá trị:} Đồ vật cất vào bên trong chiếc hộp.
\end{itemize}

\subsection*{\color{AsianBlue}\small b) Cú pháp khai báo \& Bẫy \textquotedblleft Giá trị rác\textquotedblright\ (Garbage Value)}
Trong C++, bạn bắt buộc phải khai báo kiểu dữ liệu cho biến trước khi sử dụng:
\begin{lstlisting}[language=C++]
int tuoi = 16;          // Hop chua so nguyen, gia tri 16
double diem = 8.5;      // Hop chua so thap phan
char hang = 'A';        // Hop chua 1 ky tu (trong dau nhay don)
bool da_nop = true;     // Hop chua gia tri logic: dung (true) hoac sai (false)
\end{lstlisting}

\begin{tcolorbox}[colback=red!5!white,colframe=red!75!black,arc=1mm,title=\textbf{\footnotesize BẪY NGUY HIỂM: GIÁ TRỊ RÁC (UNINITIALIZED VARIABLE)}]
\footnotesize
Nếu bạn chỉ khai báo \texttt{int tong;} mà không gán giá trị, biến sẽ chứa một giá trị rác ngẫu nhiên trong RAM (ví dụ: \texttt{-858993460}). Phép tính \texttt{tong = tong + 5;} sẽ cho kết quả sai hoàn toàn!\\
$\implies$ \textbf{Quy tắc sống còn:} Luôn khởi tạo \texttt{long long tong = 0;} và \texttt{int dem = 0;}.
\end{tcolorbox}

\subsection*{\color{AsianBlue}\small c) Các kiểu dữ liệu cốt lõi \& Bẫy tràn số nguyên (Overflow)}
\begin{center}
\small
\begin{tabular}{@{}llcp{6.5cm}@{}}
\toprule
\textbf{Kiểu} & \textbf{Kích cỡ} & \textbf{Khoảng giá trị} & \textbf{Ứng dụng trong lập trình thi đấu} \\
\midrule
\texttt{int} & 4 Bytes & $-2 \cdot 10^9 \to 2 \cdot 10^9$ & Biến đếm, chỉ số mảng, số lượng $N \le 10^6$. \\
\texttt{long long} & 8 Bytes & $-9 \cdot 10^{18} \to 9 \cdot 10^{18}$ & \textbf{Bắt buộc cho Tổng, Tích, Diện tích}. \\
\texttt{double} & 8 Bytes & Số thực chính xác cao & Điểm số, khoảng cách, diện tích hình tròn. \\
\texttt{char} & 1 Byte & 1 ký tự ASCII & Bảng chữ cái (\texttt{'A'}, \texttt{'z'}, \texttt{'0'}). \\
\texttt{bool} & 1 Byte & \texttt{true} (1) / \texttt{false} (0) & Cờ hiệu kiểm tra điều kiện. \\
\bottomrule
\end{tabular}
\end{center}

\noindent
\textbf{Mổ xẻ bẫy tràn số khi nhân:}
\begin{lstlisting}[language=C++]
int a = 1000000, b = 1000000;
long long s = a * b;          // SAI! a * b bi tran kieu int thanh so am truoc khi gan vao s!
long long dung = 1LL * a * b; // DUNG! 1LL ep toan bo phep nhan tinh tren kieu 64-bit.
\end{lstlisting}

\section*{\color{AsianDarkBlue}\normalsize 4. Các toán tử số học \& Bẫy chia nguyên}
\begin{itemize}
    \setlength{\itemsep}{2pt}
    \small
    \item \textbf{Chia nguyên (\texttt{/}) vs Chia dư (\texttt{\%}):}
    \begin{itemize}
        \item \texttt{7 / 2 = 3}: Vì cả 7 và 2 là số nguyên nên C++ vứt bỏ phần thập phân! Muốn ra số thực: \texttt{7.0 / 2 = 3.5} hoặc \texttt{(double)a / b}.
        \item \texttt{7 \% 2 = 1}: Phép chia lấy dư (chỉ dùng cho số nguyên).
    \end{itemize}
    \item \textbf{Ứng dụng kinh điển:}
    \begin{itemize}
        \item Kiểm tra chẵn lẻ: \texttt{if (n \% 2 == 0)} là số chẵn; \texttt{if (n \% 2 != 0)} là số lẻ.
        \item Tách chữ số hàng đơn vị: \texttt{n \% 10} lấy chữ số cuối; \texttt{n / 10} xóa bỏ chữ số cuối.
    \end{itemize}
    \item \textbf{Phép toán rút gọn:} \texttt{x++} (tăng 1), \texttt{x--} (giảm 1), \texttt{tong += a} (cộng dồn), \texttt{tich *= a}.
\end{itemize}

\newpage

\section*{\color{AsianDarkBlue}\normalsize 5. Nhập / Xuất chuẩn mực trong C++}
\begin{itemize}
    \setlength{\itemsep}{2pt}
    \small
    \item \textbf{Xuất dữ liệu (\texttt{cout <<}):} Dấu \texttt{<<} đẩy dữ liệu ra màn hình. Dùng \texttt{\textbackslash n} để xuống dòng (nhanh hơn \texttt{endl}).
    \item \textbf{Nhập dữ liệu (\texttt{cin >>}):} Dấu \texttt{>>} hút dữ liệu từ bàn phím vào biến, tự động bỏ qua dấu cách và dấu xuống dòng. Nhập nhiều biến: \texttt{cin >> a >> b >> c;}.
    \item \textbf{In số thập phân chính xác:} Dùng \texttt{<iomanip>}: \texttt{cout << fixed << setprecision(2) << x;}.
\end{itemize}

\section*{\color{AsianDarkBlue}\normalsize 6. Cấu trúc rẽ nhánh \texttt{if / else} \& Tư duy Logic}

\subsection*{\color{AsianBlue}\small a) Toán tử so sánh \& Phép toán logic}
\begin{itemize}
    \setlength{\itemsep}{1.5pt}
    \small
    \item So sánh: \texttt{==} (so sánh bằng, \textbf{tuyệt đối không viết nhầm thành 1 dấu =}), \texttt{!=} (khác), \texttt{>}, \texttt{>=}, \texttt{<}, \texttt{<=}.
    \item Logic: \texttt{\&\&} (VÀ -- đồng thời đúng), \texttt{||} (HOẶC -- ít nhất một điều kiện đúng), \texttt{!} (PHỦ ĐỊNH).
\end{itemize}

\begin{tcolorbox}[colback=red!5!white,colframe=red!75!black,arc=1mm,title=\textbf{\footnotesize BẪY TOÁN HỌC TRONG C++: ĐIỀU KIỆN KẸP GIỮA}]
\footnotesize
Trong toán ta viết: $1 \le x \le 10$. Trong C++, viết \texttt{if (1 <= x <= 10)} là \textbf{HOÀN TOÀN SAI}!\\
(Máy tính tính \texttt{1 <= x} ra 0 hoặc 1, sau đó so sánh tiếp với 10 $\implies$ luôn đúng).\\
$\implies$ \textbf{Cách viết chuẩn:} \texttt{if (1 <= x \&\& x <= 10)}.
\end{tcolorbox}

\subsection*{\color{AsianBlue}\small b) Khung rẽ nhánh mẫu}
\begin{lstlisting}[language=C++]
if (diem >= 8.0) {
    cout << "GIOI\n";
} else if (diem >= 6.5) {
    cout << "KHA\n";
} else if (diem >= 5.0) {
    cout << "TRUNG BINH\n";
} else {
    cout << "YEU\n";
}
\end{lstlisting}

\section*{\color{AsianDarkBlue}\normalsize 7. Vòng lặp \texttt{for} \& Kỹ thuật Chạy từng bước bằng tay (Dry Run)}
\subsection*{\color{AsianBlue}\small a) Cơ chế hoạt động của vòng \texttt{for}}
Cú pháp: \texttt{for (Khởi\_tạo ; Điều\_kiện\_lặp ; Bước\_nhảy) \{ Thân vòng lặp \}}
\begin{enumerate}
    \setlength{\itemsep}{1pt}
    \small
    \item Thực hiện \texttt{Khởi\_tạo} duy nhất 1 lần ở đầu.
    \item Kiểm tra \texttt{Điều\_kiện\_lặp}: Đúng $\to$ thực hiện thân vòng lặp. Sai $\to$ thoát ngay.
    \item Chạy xong thân vòng lặp $\to$ thực hiện \texttt{Bước\_nhảy} (\texttt{i++}) $\to$ quay lại kiểm tra điều kiện.
\end{enumerate}

\subsection*{\color{AsianBlue}\small b) Bảng Trace Dry Run: Tính tổng $S = 1 + 2 + 3 + 4$}
\begin{lstlisting}[language=C++]
long long S = 0;
for (int i = 1; i <= 4; i++) {
    S += i;
}
cout << S << "\n";
\end{lstlisting}

\begin{center}
\footnotesize
\begin{tabular}{@{}ccccccc@{}}
\toprule
\textbf{Vòng} & \textbf{Biến $i$} & \textbf{Điều kiện $i \le 4$} & \textbf{$S$ trước cộng} & \textbf{Lệnh $S = S + i$} & \textbf{$S$ sau cộng} & \textbf{Bước nhảy $i$++} \\
\midrule
Khởi tạo & - & - & - & $S = 0$ & 0 & $i = 1$ \\
1 & 1 & $1 \le 4$ (\textbf{Đúng}) & 0 & $S = 0 + 1$ & \textbf{1} & $i$ tăng lên 2 \\
2 & 2 & $2 \le 4$ (\textbf{Đúng}) & 1 & $S = 1 + 2$ & \textbf{3} & $i$ tăng lên 3 \\
3 & 3 & $3 \le 4$ (\textbf{Đúng}) & 3 & $S = 3 + 3$ & \textbf{6} & $i$ tăng lên 4 \\
4 & 4 & $4 \le 4$ (\textbf{Đúng}) & 6 & $S = 6 + 4$ & \textbf{10} & $i$ tăng lên 5 \\
Thoát & 5 & $5 \le 4$ (\textbf{Sai!}) & 10 & \textit{Không chạy} & \textbf{10} & Dừng vòng lặp \\
\bottomrule
\end{tabular}
\end{center}

\newpage

\section*{\color{AsianDarkBlue}\normalsize 8. Hướng dẫn Kỹ thuật Đọc/Ghi Tệp thi đấu: Từ Cơ bản đến Nâng cao (Tệp \texttt{.INP} / \texttt{.OUT})}

Trong kỳ thi Olympic Tin học, quy định bắt buộc 100\% các bài thi phải đọc dữ liệu từ tệp \texttt{TEN\_BAI.INP} và xuất ra tệp \texttt{TEN\_BAI.OUT}. Nếu in ra màn hình hoặc đọc sai tên tệp, bài thi sẽ nhận \textbf{0 điểm}.

\vspace{-1mm}
\subsection*{\color{AsianBlue}\footnotesize Cấp độ 1 (Cơ bản): Chuyển hướng luồng dữ liệu bằng \texttt{freopen}}
Cơ chế của hàm \texttt{freopen} trong \texttt{<cstdio>} là \textbf{bẻ hướng dòng dữ liệu}:
\begin{itemize}
    \setlength{\itemsep}{0pt}
    \footnotesize
    \item \texttt{freopen("TONG.INP", "r", stdin);} $\implies$ Đọc từ file thay vì bàn phím.
    \item \texttt{freopen("TONG.OUT", "w", stdout);} $\implies$ Ghi vào file thay vì màn hình. Lệnh \texttt{cin}, \texttt{cout} \textbf{giữ nguyên 100\%}.
\end{itemize}

\vspace{-2mm}
\subsection*{\color{AsianBlue}\footnotesize Cấp độ 2 (Trung cấp): Tự động nhận diện tệp thông minh bằng \texttt{if (fopen)}}
Dùng \texttt{if (fopen("TONG.INP", "r"))} trước \texttt{freopen}. Có file $\to$ đọc file; không có file $\to$ nhập bàn phím. Học sinh \textbf{không cần sửa hay đóng/mở comment} khi nộp thi hay test ở nhà.

\vspace{-2mm}
\subsection*{\color{AsianBlue}\footnotesize Cấp độ 3 (Nâng cao): Đọc đến hết tệp (EOF) \& Tối ưu tốc độ Fast I/O}
\begin{itemize}
    \setlength{\itemsep}{0pt}
    \footnotesize
    \item \textbf{Đọc hết tệp (EOF):} Khi không cho trước số phần tử $N$, dùng vòng lặp: \texttt{while (cin >> x) \{ ... \}}.
    \item \textbf{Tối ưu Fast I/O:} Thêm \texttt{ios\_base::sync\_with\_stdio(false); cin.tie(NULL);} và dùng \texttt{"\textbackslash n"} (tránh TLE).
\end{itemize}

\vspace{-1mm}
\noindent
\begin{minipage}[t]{0.48\textwidth}
\textbf{\color{AsianDarkBlue}\footnotesize Code mẫu 1: Cấp độ 1 (Cơ bản cố định)}
\begin{lstlisting}[language=C++,basicstyle=\ttfamily\scriptsize,numbers=none,aboveskip=1mm,belowskip=1mm]
#include <iostream>
#include <cstdio>
using namespace std;
int main() {
    freopen("TONG.INP", "r", stdin);
    freopen("TONG.OUT", "w", stdout);
    long long a, b;
    if (cin >> a >> b) {
        cout << a + b << "\n";
    }
    return 0;
}
\end{lstlisting}
\end{minipage}
\hfill
\begin{minipage}[t]{0.49\textwidth}
\textbf{\color{AsianDarkBlue}\footnotesize Code mẫu 2: Chuẩn Olympic (Tự động + EOF)}
\begin{lstlisting}[language=C++,basicstyle=\ttfamily\scriptsize,numbers=none,aboveskip=1mm,belowskip=1mm]
#include <iostream>
#include <cstdio>
using namespace std;
int main() {
    ios_base::sync_with_stdio(false);
    cin.tie(NULL);
    if (fopen("TONG.INP", "r")) {
        freopen("TONG.INP", "r", stdin);
        freopen("TONG.OUT", "w", stdout);
    }
    long long a, tong = 0;
    while (cin >> a) tong += a;
    cout << tong << "\n";
    return 0;
}
\end{lstlisting}
\end{minipage}

\vspace{-2mm}
\section*{\color{AsianDarkBlue}\normalsize 9. Bảng chẩn đoán lỗi thường gặp trên Code::Blocks cho học sinh mới}
\begin{center}
\footnotesize
\begin{tabular}{@{}p{4.2cm}p{4.5cm}p{7.3cm}@{}}
\toprule
\textbf{Thông báo lỗi trên màn hình} & \textbf{Nguyên nhân thực tế} & \textbf{Cách xử lý ngay lập tức} \\
\midrule
\texttt{error: expected ';' before...} & Quên dấu \texttt{;} ở cuối dòng liền trước. & Nhìn lên dòng phía trên và thêm \texttt{;}. \\
\addlinespace[2pt]
\texttt{error: 'x' was not declared...} & Biến \texttt{x} chưa khai báo, hoặc sai hoa/thường. & Khai báo kiểu (\texttt{int x;}) trước khi dùng. \\
\addlinespace[2pt]
\texttt{error: redefinition of 'int x'} & Khai báo biến \texttt{x} hai lần trong một hàm. & Bỏ chữ \texttt{int} ở lần gán thứ hai. \\
\addlinespace[2pt]
Kết quả ra số âm (ví dụ \texttt{-2147483648}) & Tràn số nguyên kiểu \texttt{int} ($> 2 \cdot 10^9$). & Đổi sang kiểu \texttt{long long}, ép kiểu \texttt{1LL * a * b}. \\
\addlinespace[2pt]
Cửa sổ đen console treo đứng & Vòng lặp vô tận (quên \texttt{i++} hoặc sai điều kiện). & Bấm \texttt{Ctrl + C} tắt cửa sổ, sửa biến lặp. \\
\bottomrule
\end{tabular}
\end{center}

\end{document}
